\section{The weakly universal variant}
\label{sec:weak}

The strongly universal automaton of Section~\ref{sec:strong} contains a simpler one. Drop the track-building
rule (B1) and (B2), keep only the motion and switch rules of Section~\ref{sec:automaton}, and supply the
registers as infinite track. The result is weakly universal, and it is worth stating on its own, because it is
the cleaner object when a finite seed is not needed and because it is the direct counterpart of Margenstern's
constructions.

\subsection{Its structure}

The weak automaton has the dynamic alphabet of Table~\ref{tab:states}, empty, clear, head, trailing, and the
switch selection bit, and exactly the rules of Tables~\ref{tab:motion} and~\ref{tab:switch}. There is no build
rule, so empty space stays empty for all time and the circuit is whatever the initial configuration laid down.
A register is an infinite chain of flip-flop bits, written out in advance as an ultimately-periodic pattern, the
fixed program repeated outward to give unbounded memory. Nothing is ever constructed.

\subsection{Its theorem}

\begin{theorem}[Uniform weak universality]
\label{thm:weak}
The automaton with the rules of Tables~\ref{tab:motion} and~\ref{tab:switch} alone is weakly universal on every
regular hyperbolic tiling. For each tiling and each register machine there is an infinite, ultimately-periodic
initial configuration on that tiling whose evolution simulates the machine.
\end{theorem}

By Proposition~\ref{prop:assembly} the motion and switch rules realize the registers and the sequencer. The
registers, given as infinite ultimately-periodic track, supply the unbounded memory, so the configuration is
infinite but ultimately periodic and the universality is weak \cite{minsky1967}. As before, no rule mentions the
tiling. \qed

\subsection{Why keep both}

The weak variant is smaller, four dynamic symbols and the switch bit, with no build rule, and it is the honest
counterpart to compare against the rotation-invariant per-tiling automata, which are also weakly universal. The
strong variant adds one rule, the builder, and in exchange starts from a finite seed. They share the same
locomotive and the same three switches. One can read the weak automaton as the strong automaton restricted to
configurations that never need to grow.
